%
%
%
\section{Exercises}
\label{chap:exercises}

The following exercises are intended to consolidate the main ideas developed in
these notes.  Exercises marked \textbf{Pen-and-paper} can be completed without
specialized software.  Those marked \textbf{Simple script} require only basic
array operations, linear algebra and Fourier transforms, for example as
provided by NumPy.  When an exercise carries both labels, it is useful to make
a prediction by hand before checking it numerically.

%
%
%
\subsection{Plane waves and reciprocal space}

\begin{exci}{Reciprocal lattice of a non-cubic cell \hfill\textnormal{\small Pen-and-paper}}
    \item Consider the unit-cell matrix
    \begin{equation*}
        \mat{M}=\begin{pmatrix}
            4&0&0\\
            1&3&0\\
            0&1&5
        \end{pmatrix}.
    \end{equation*}
    Calculate the unit-cell volume $\Omega$ and the reciprocal-lattice matrix
    $\mat{B}$.
    \item Verify explicitly that
    $\vec{a}_{i}\cdot\vec{b}_{j}=2\pi\delta_{ij}$.
\end{exci}

\begin{exci}{Enumerating FFT modes \hfill\textnormal{\small Pen-and-paper}}
    \item List the integer frequency indices for $N=5$, $N=6$ and $N=8$
    before and after shifting the zero-frequency component to the centre.
    \item Identify the Nyquist mode for each even grid and explain why the
    positive and negative Nyquist frequencies constitute only one independent
    mode on the sampling grid.
\end{exci}

\begin{exci}{Orthonormality in fractional coordinates \hfill\textnormal{\small Pen-and-paper}}
    \item Write a point in an arbitrary parallelepiped unit cell as
    $\vec{r}=s_{1}\vec{a}_{1}+s_{2}\vec{a}_{2}+s_{3}\vec{a}_{3}$, with
    $0\leq s_i<1$.
    \item Use these fractional coordinates, including the appropriate
    Jacobian, to prove the plane-wave orthonormality relation
    \begin{equation*}
        \braket{\phi(\vec{G}\prime)|\phi(\vec{G})}
        =\delta_{\vec{G},\vec{G}\prime}.
    \end{equation*}
\end{exci}

\begin{exci}{Bloch phases \hfill\textnormal{\small Pen-and-paper}}
    \item For a one-dimensional lattice with spacing $a$, calculate the phase
    acquired under translations by $a$, $2a$ and $3a$ for
    \begin{equation*}
        k=0,\qquad k=\frac{\pi}{2a},\qquad k=\frac{\pi}{a}.
    \end{equation*}
    \item Show that $k$ and $k+2\pi/a$ impose the same boundary condition.
\end{exci}

\begin{exci}{Counting plane waves \hfill\textnormal{\small Pen-and-paper / Simple script}}
    \item For a cubic unit cell, enumerate all reciprocal vectors satisfying
    $G^{2}/2<E_{\textrm{cut}}$ for several cutoff energies.
    \item Compare the number of vectors with the continuum estimate
    \begin{equation*}
        N_{\textrm{PW}}\approx
        \frac{\Omega}{6\pi^{2}}(2E_{\textrm{cut}})^{3/2}.
    \end{equation*}
    \item Explain the small irregularities, or ``shell effects'', in the exact
    counts.
\end{exci}

\begin{exci}{Aliasing by inspection \hfill\textnormal{\small Pen-and-paper / Simple script}}
    \item Sample $f_m(x)=\exp(2\pi i m x/L)$ at $N$ equally spaced points in a
    periodic interval of length $L$.
    \item Show algebraically that modes $m$ and $m+N$ give identical samples.
    \item Verify the result for several choices of $m$ and $N$ using a short
    script.
\end{exci}

\begin{exci}{FFT normalization and Parseval's theorem \hfill\textnormal{\small Simple script}}
    \item Sample the periodic function
    \begin{equation*}
        f(x)=1+2\cos(2\pi x/L)-\sin(4\pi x/L)
    \end{equation*}
    on an equidistant grid and transform it with \texttt{numpy.fft.fft} or an
    equivalent FFT implementation.
    \item Identify all non-zero Fourier coefficients and reconstruct $f(x)$.
    \item Determine the normalization needed to convert the library output to
    the plane-wave convention used in these notes.
    \item Verify Parseval's identity numerically.
\end{exci}

\begin{exci}{Translation becomes phase \hfill\textnormal{\small Pen-and-paper / Simple script}}
    \item Translate the function in the preceding exercise by a displacement
    $d$ and predict the change in each Fourier coefficient.
    \item Verify the prediction numerically and confirm that the magnitudes of
    the coefficients remain unchanged.
\end{exci}

\begin{exci}{Why the density needs a finer grid \hfill\textnormal{\small Pen-and-paper / Simple script}}
    \item Let
    \begin{equation*}
        \psi(x)=\cos(G_{\max}x)
        +\frac{1}{2}\cos\big[(G_{\max}-\Delta G)x\big].
    \end{equation*}
    Expand $|\psi(x)|^{2}$ by hand and identify its largest frequency.
    \item Evaluate the density numerically on the original orbital grid and
    demonstrate the aliasing that occurs when the grid cannot represent this
    largest frequency.
\end{exci}

%
%
%
\subsection{Potentials and electrostatics}

\begin{exci}{Poisson's equation for one Fourier mode \hfill\textnormal{\small Pen-and-paper}}
    \item Given
    $\rho(\vec{r})=\rho_{0}\cos(\vec{G}_{0}\cdot\vec{r})$, solve Poisson's
    equation analytically for the electrostatic potential.
    \item Verify the result by applying the Laplacian and explain why this
    example has no problematic $\vec{G}=\vec{0}$ contribution.
\end{exci}

\begin{exci}{Meaning of the zero mode \hfill\textnormal{\small Pen-and-paper}}
    \item Show that adding a constant to an electrostatic potential does not
    change its electric field.
    \item Explain the distinction between choosing the average potential to be
    zero, removing the average charge density, and requiring a periodically
    repeated unit cell to be neutral.
\end{exci}

\begin{exci}{Hartree energy in reciprocal space \hfill\textnormal{\small Pen-and-paper}}
    \item Starting from
    \begin{equation*}
        E_{\textrm{H}}=\frac{1}{2}\int_{\Omega}
        \rho(\vec{r})V_{\textrm{H}}(\vec{r})\,d\vec{r},
    \end{equation*}
    derive an expression for $E_{\textrm{H}}$ in terms of the reciprocal-space
    density coefficients.
    \item Use the result to show that the Hartree energy of a neutral,
    real-valued density fluctuation is non-negative.
\end{exci}

\begin{exci}{Electrostatic cancellation \hfill\textnormal{\small Simple script}}
    \item Place a normalized Gaussian electron density and an oppositely
    charged ion in a cubic unit cell.
    \item Calculate the electron--electron and electron--ion contributions
    separately as the cell size changes, and plot their sum.
    \item Explain why the separate contributions converge more poorly than the
    neutral combination.
\end{exci}

\begin{exci}{Functional derivative of LDA exchange \hfill\textnormal{\small Pen-and-paper}}
    \item Starting from
    \begin{equation*}
        E_{x}[\rho]=-C_{x}\int \rho^{4/3}(\vec{r})\,d\vec{r},
    \end{equation*}
    calculate $\delta E_x/\delta\rho(\vec{r})$.
    \item Determine its sign and explain its physical meaning.
    \item Compare your result carefully with the exchange-potential convention
    used in the main text.
\end{exci}

%
%
%
\subsection{GTH pseudopotentials}

\begin{exci}{Limits of the local GTH potential \hfill\textnormal{\small Pen-and-paper}}
    \item Use the small-argument expansion of the error function to calculate
    $V_{\textrm{loc}}(0)$.
    \item Determine the limit as $r\rightarrow\infty$ and identify which
    polynomial--Gaussian terms survive at the origin.
    \item Explain how these two limits express the intended purpose of the
    pseudopotential.
\end{exci}

\begin{exci}{The finite $G=0$ contribution \hfill\textnormal{\small Pen-and-paper}}
    \item Expand the reciprocal-space local GTH potential around $G=0$.
    \item Separate its divergent Coulomb contribution from its finite
    remainder and recover the finite constant stated in the text.
\end{exci}

\begin{exci}{A two-atom structure factor \hfill\textnormal{\small Pen-and-paper / Simple script}}
    \item For identical atoms at $\pm\vec{R}/2$, simplify
    \begin{equation*}
        S(\vec{G})=
        \exp(-i\vec{G}\cdot\vec{R}/2)
        +\exp(+i\vec{G}\cdot\vec{R}/2).
    \end{equation*}
    \item Determine which reciprocal components vanish for selected atomic
    separations.
    \item Explain why the corresponding real-space potential is real and
    verify this with a small numerical transform.
\end{exci}

\begin{exci}{Non-local projectors as a low-rank operator \hfill\textnormal{\small Pen-and-paper}}
    \item Choose two projector vectors, collect them as the columns of $B$,
    and choose a real symmetric $2\times2$ matrix $h$. Apply $BhB^{\dagger}$
    to a trial vector.
    \item Prove that $BhB^{\dagger}$ is Hermitian and has rank no greater than
    two.
    \item Express its expectation value without constructing the full
    $N_{\textrm{PW}}\times N_{\textrm{PW}}$ matrix.
\end{exci}

\begin{exci}{Reading a pseudopotential file \hfill\textnormal{\small Pen-and-paper}}
    \item For the carbon GTH parameter file shown in the text, identify
    $Z_{\textrm{ion}}$, $r_{\textrm{loc}}$, all active $C_n$, the angular
    channels, the number of projectors and every active $h^l$ matrix.
    \item How many electrons and doubly occupied spatial orbitals would be
    treated explicitly in a closed-shell calculation of \ce{CH4} using these
    pseudopotentials?
\end{exci}

%
%
%
\subsection{Ewald summation}

\begin{exci}{Gaussian screening charge \hfill\textnormal{\small Pen-and-paper}}
    \item Normalize the three-dimensional Gaussian screening charge used in
    the Ewald construction.
    \item Derive its Fourier transform and check its dimensions and its limits
    as $G\rightarrow0$ and $G\rightarrow\infty$.
\end{exci}

\begin{exci}{Deriving the complementary-error-function split \hfill\textnormal{\small Pen-and-paper}}
    \item Starting from the potential of a Gaussian charge, show that
    \begin{equation*}
        \frac{1}{r}=
        \frac{\erfc(\sqrt{\alpha}r)}{r}
        +\frac{\erf(\sqrt{\alpha}r)}{r}.
    \end{equation*}
    \item Explain which term is most efficiently evaluated in real space and
    which in reciprocal space.
\end{exci}

\begin{exci}{Enumerating a primed Ewald sum \hfill\textnormal{\small Pen-and-paper}}
    \item Consider two ions in the central cell and one neighboring cell on
    either side in a one-dimensional sketch. Explicitly list the interactions
    within the central cell, those with periodic images and the excluded
    self-interactions.
    \item Write the result as a compact primed summation and explain precisely
    what the prime excludes.
\end{exci}

\begin{exci}{Independence of the splitting parameter \hfill\textnormal{\small Simple script}}
    \item Implement the real-space, reciprocal-space, self-interaction and
    background contributions for two opposite charges in a cubic cell.
    \item Vary $\alpha$ while keeping both truncation ranges sufficiently
    large. Plot the separate contributions and their sum.
    \item Explain why the individual terms change while the total energy
    remains approximately constant.
\end{exci}

\begin{exci}{Balancing real- and reciprocal-space cost \hfill\textnormal{\small Pen-and-paper / Simple script}}
    \item Use the parameter-estimation formulas in the text to tabulate the
    real-space and reciprocal-space ranges for several splitting parameters.
    \item Estimate the number of terms in each sum and identify the
    approximately cheapest choice.
\end{exci}

%
%
%
\subsection{Eigensolvers and self-consistency}

\begin{exci}{Two Arnoldi steps by hand \hfill\textnormal{\small Pen-and-paper}}
    \item Choose a small Hermitian $3\times3$ matrix and a normalized starting
    vector. Perform two Arnoldi iterations and construct $Q_2$ and $H_2$.
    \item Calculate the Ritz pairs of $H_2$ and form the corresponding Ritz
    vectors in the original space.
    \item Evaluate one Ritz residual both directly and from the Arnoldi
    residual formula.
\end{exci}

\begin{exci}{Planning a matrix-free Hamiltonian application \hfill\textnormal{\small Pen-and-paper}}
    \item For each of the kinetic, local ionic, Hartree,
    exchange--correlation and non-local projector contributions, state the
    representation in which it is cheapest to apply.
    \item Describe the operation that must be performed and identify every
    point at which an FFT or inverse FFT is required.
    \item Combine the operations into pseudocode for one Hamiltonian--vector
    product.
\end{exci}

\begin{exci}{Stability of linear mixing \hfill\textnormal{\small Pen-and-paper}}
    \item Replace the density by a scalar and let the SCF map be
    $\rho_{\textrm{out}}=a\rho_{\textrm{in}}+b$. Derive its fixed point.
    \item Show that linear mixing converges when
    \begin{equation*}
        \left|1+\beta(a-1)\right|<1.
    \end{equation*}
    \item Construct an example for which $\beta=1$ diverges but $\beta=0.2$
    converges.
\end{exci}

\begin{exci}{Pulay coefficients for two residuals \hfill\textnormal{\small Pen-and-paper}}
    \item Choose two short residual vectors, construct their overlap matrix
    $B$, and solve the constrained Pulay system.
    \item Compare the norm of the combined residual with the two original
    norms.
    \item Repeat the calculation for nearly parallel residuals and explain the
    resulting ill-conditioning.
\end{exci}

\begin{exci}{Charge conservation during mixing \hfill\textnormal{\small Pen-and-paper / Simple script}}
    \item Apply the density-rescaling formula to a small grid of trial density
    values and verify the electron count after rescaling.
    \item Construct an example showing why normalization alone cannot repair a
    candidate density containing negative values.
\end{exci}

\begin{exci}{Phase rotation without an optimizer \hfill\textnormal{\small Pen-and-paper / Simple script}}
    \item Write $\psi_k=x_k+iy_k$ and show that
    \begin{align*}
        \sum_k\left[\mathfrak{R}\left(e^{i\varphi}\psi_k\right)\right]^2
        ={}&\frac{A+B}{2}+\frac{A-B}{2}\cos(2\varphi)\\
        &-C\sin(2\varphi),
    \end{align*}
    where $A=\sum_kx_k^2$, $B=\sum_ky_k^2$ and $C=\sum_kx_ky_k$.
    \item Obtain the optimal phase analytically and compare it with the result
    of a direct numerical search.
    \item Explain why a general nonlinear optimizer is unnecessary for this
    one-parameter problem.
\end{exci}

\begin{exci}{Fourier zero-padding \hfill\textnormal{\small Simple script}}
    \item Upsample a periodic function by padding its reciprocal-space
    coefficients with zeros and applying an inverse FFT on the enlarged grid.
    \item Verify that the original values are recovered at matching grid
    points and determine the normalization required to conserve the norm.
    \item Explain why zero-padding interpolates a band-limited function but
    adds no new physical information.
\end{exci}
